\documentclass[11pt,a4paper]{article}
\usepackage[utf8]{inputenc}
\usepackage[T1]{fontenc}
\usepackage[italian]{babel}
\usepackage{amsmath,amssymb,graphicx,booktabs,siunitx,hyperref}
\usepackage{listings}
\lstdefinestyle{faust}{basicstyle=\ttfamily\small,breaklines=true,
  keepspaces=true,columns=fullflexible,frame=single,framesep=4pt}
\lstset{style=faust}
\usepackage[margin=2.5cm]{geometry}
\graphicspath{{figures/}}

\title{dslar: l'ecosistema di Larsen autoregolato di Di Scipio\\\large Diario di studio e ricostruzione per la suite SEAM-LTM}
\author{Giuseppe Silvi --- SEAM}
\date{2026}

\begin{document}
\maketitle
\begin{abstract}
Questo documento è un diario di studio passo passo.
Analizza la patch \texttt{LAR.pd} di Agostino Di Scipio a partire dal suo sorgente e dalla letteratura primaria dell'autore.
Ne ricostruisce il motore di autoregolazione del feedback acustico, la legge di controllo omeostatica e la strategia di indipendenza dalla frequenza di campionamento, preparando le fasi successive: libreria Faust \texttt{sds}, esempio compilabile e plugin.
\end{abstract}
\tableofcontents

\section{Introduzione e motivazione}\label{sec:intro}
Segnaposto: sostituito in Task 2.

\section{Anatomia della patch}\label{sec:anatomia}
Segnaposto: scritto in Task 2.

\section{Contesto e fonti}\label{sec:contesto}
Segnaposto: scritto in Task 3.

\section{Analisi del segnale}\label{sec:segnale}
Segnaposto: scritto in Task 4.

\section{La legge di controllo omeostatica}\label{sec:controllo}
Segnaposto: scritto in Task 4.

\section{Compatibilità con la frequenza di campionamento}\label{sec:sr}
Segnaposto: scritto in Task 4.

\section{Mappa di decomposizione}\label{sec:decomposizione}
Segnaposto: scritto in Task 5.

\section{Sintesi e ponte alla Fase 2}\label{sec:sintesi}
Segnaposto: scritto in Task 5.

\bibliographystyle{plain}
\bibliography{refs}
\end{document}
